\chapter{Methodology}
\label{chap:methodology}

This chapter explains the methods used in this study. It describes how
the evaluation dataset was built, which Large Language Models were chosen
and why, how the Chain-of-Thought and Tree-of-Thought prompting
strategies were designed, how answer correctness was judged, which
statistical tests were used for each research question, how the
error-judging process was set up, and the limitations of these methods.

\section{Dataset Construction}
\label{sec:dataset}

The evaluation dataset includes 131 mathematical problems, grouped into
three difficulty levels (Easy, Medium, Hard) and four categories: Algebra,
Number Theory, Probability, and Combinatorics. Table~\ref{tab:composition}
lists the exact breakdown.

\begin{table}[H]
  \centering
  \caption{Dataset composition by category and difficulty level.}
  \label{tab:composition}
  \begin{tabular}{lrrrr}
    \toprule
    Category & Easy & Medium & Hard (AIME) & Total \\
    \midrule
    Algebra          & 10 & 10 & 16 & 36 \\
    Number Theory    & 10 & 10 & 11 & 31 \\
    Probability      & 10 & 10 & 8  & 28 \\
    Combinatorics    & 10 & 10 & 16 & 36 \\
    \midrule
    \textbf{Total}   & \textbf{40} & \textbf{40} & \textbf{51} & \textbf{131} \\
    \bottomrule
  \end{tabular}
\end{table}

\begin{figure}[H]
  \centering
  \includegraphics[width=0.7\textwidth]{dataset_composition_pie.png}
  \caption{Dataset composition by difficulty tier and category.}
  \label{fig:dataset-composition-pie-methodology}
\end{figure}

The Easy and Medium levels come from Hendrycks MATH, with exactly 10
questions per category in each level (using seed 25116096). These match
Hendrycks Levels 1--2 for Easy and Levels 3--4 for Medium. The Hard level
uses only questions from AIME.

The Hard level uses AIME questions only, rather than also including
Hendrycks MATH Level 5. AIME is a single, consistent competition-mathematics
source with a uniform integer answer format (0--999), which keeps
answer-parsing and grading consistent across the whole Hard tier. Hendrycks
MATH Level 5 is a separate item bank with a different answer format
(including fractional answers) and no independent validation that its
internal ``Level 5'' label represents the same difficulty as an AIME
problem. Mixing the two sources under one difficulty label would risk
treating non-comparable item banks as equivalent, so the Hard level was
restricted to AIME only.

The 51 Hard-level questions were drawn from six AIME contests: 2024 I,
2024 II, 2025 I, 2025 II, 2026 I, and 2026 II. The number of questions in
each Hard-level category is intentionally uneven, reflecting how many
questions of each category were available and labelled across these six
contests. Probability is limited to 8 questions because that is all the
AIME Probability questions available and labelled for this project. The
other three categories were not reduced to match. Reducing them would
have meant losing useful data without any real benefit, so the imbalance
is handled as a covariate in the analysis instead of forcing equal
numbers.

For most analyses in this study, the Hard tier is further restricted to
a fixed, seeded random sample of 40 of its 51 questions (seed 25116096),
matching the fixed size of the Easy and Medium tiers (40 questions each).
This gives a consistent, equal-sized pool of 40 questions per difficulty
tier -- 120 questions in total, so that no single tier's larger
question count disproportionately influences comparisons across tiers.
This balanced 120-question pool is used throughout this dissertation's
analyses unless otherwise stated. The one deliberate exception is the
per-question correctness matrix (Section~\ref{sec:question-correctness}),
which uses the full 51 Hard-tier questions (131 questions in total) to
show every curated question rather than a sample of them.

\section{Models}
\label{sec:models}

This study evaluates three large language models: Gemini Flash Lite
(\texttt{gemini-flash-lite-latest}, resolving to
\texttt{gemini-3.5-flash-lite}), GPT-OSS-120B (\texttt{openai/gpt-oss-120b},
via Groq), and Mistral Large (\texttt{mistral-large-latest}, resolving to
Mistral Large 3, \texttt{mistral-large-2512}). Each model is accessed via
its provider's free-tier API.

Gemini was included to ensure diversity among model providers. Groq hosts
open-weight models, originally from Meta's Llama family and later
OpenAI's GPT-OSS, while Mistral AI offers its proprietary Mistral-family
models. Including Gemini introduces a model from Google's family,
avoiding reliance on adjacent open-weight or single-vendor lineages. The
specific Groq and Mistral models were chosen after an earlier review
determined that the initially planned Llama-3.3-70B (Groq) and Mistral
Small performed inadequately on AIME-level problems. Additionally,
Llama-4-Maverick/Scout, the intended replacement on Groq, was deprecated.
GPT-OSS-120B was selected as Groq's recommended migration target for
these deprecated models.

All models are run with temperature set to 0, ensuring deterministic
responses for identical prompts. Two of the three models, Gemini and
GPT-OSS-120B, provide an internal control for ``thinking effort'' or
``reasoning effort.'' Both are set to their lowest available settings
(Gemini: \texttt{thinking\_level = "minimal"}; GPT-OSS-120B:
\texttt{reasoning\_effort = "low"}) and remain constant across both
Chain-of-Thought and Tree-of-Thought conditions. This control is
particularly relevant for RQ3. If this setting were allowed to vary,
differences in accuracy between Chain-of-Thought and Tree-of-Thought
could result from changes in internal reasoning budget rather than the
prompting strategy. Fixing the setting at its minimum value controls for
this confound, though it does not eliminate it entirely. This approach
does not guarantee the absence of internal reasoning, but ensures that
any such reasoning remains consistent across conditions. Mistral Large
does not expose an equivalent control and is used as a standard instruct
model throughout.

\section{Prompting Strategies}
\label{sec:prompting}
\subsection{Chain-of-Thought}

In the Chain-of-Thought setup, the model answers each question in a
single call. The model receives the problem and is instructed to explain
its reasoning step by step, concluding with a final answer in a
specified format.

This approach follows the standard single-pass Chain-of-Thought design,
in which one call produces a complete reasoning process and a final
answer. All outputs are generated at temperature 0 to guarantee
consistency. The prompt template standardises the output format,
enabling uniform parsing of responses among different models:

\begin{quote}
\ttfamily\small
Solve the following mathematics problem.\\
\\
Show your reasoning as clear, numbered steps.\\
\\
Format your response exactly as:\\
\\
Step 1: ...\\
Step 2: ...\\
Step 3: ...\\
\\
Final Answer: <answer>\\
\\
Problem:\\
\{question\}
\end{quote}

\subsection{Tree-of-Thought Design}
\label{sec:tot-design}

In the Tree-of-Thought setup, the model generates three independent
candidate solutions, referred to as branches. Subsequently, an
additional call selects the optimal branch. Thus, four calls are
executed per question, compared with the single call used in
Chain-of-Thought (CoT) prompting.

For each of the three branches, the model uses the same step-by-step
prompt as in CoT, with an additional instruction: each response is
explicitly labelled as a separate solution attempt (e.g., ``this is one
independent solution attempt, attempt N of 3''). The model is instructed
to commit to a single approach for each attempt, avoiding hedging
between methods. Branches are generated at a temperature of 0.7, rather
than 0, to promote diversity among solutions. A temperature of 0 would
result in nearly identical outputs, reducing the effectiveness of
comparative evaluation. Employing a moderate temperature ensures that
the branches represent genuinely distinct solutions.

After generating all three branches, the model performs an additional
call at temperature 0 to maintain consistency. During this step, the
model reviews the complete text of all three candidate solutions and
evaluates which branch demonstrates the most mathematically sound
reasoning. The model is instructed to select the branch with the
strongest justification, rather than defaulting to the most frequently
occurring answer in cases of disagreement among branches:

\begin{quote}
\ttfamily\small
Evaluate each candidate branch for mathematical soundness. Select the
ONE branch most likely to be correct and well-reasoned. If multiple
branches reach the same final answer through valid reasoning, prefer
the clearest derivation. If branches disagree, select based on which
reasoning is most rigorous and free of errors --- do not simply pick
the majority answer without checking the reasoning.
\end{quote}

It returns the number of the chosen branch and a brief explanation
(``Selected Branch: \texttt{<branch number>}'', ``Justification:
\texttt{<one sentence, maximum 25 words>}''). The full text of the
selected branch is then used just like a CoT response. If this selection
response cannot be parsed into a branch number, for example if the model
does not follow the requested output format, the pipeline falls back to
using the first generated branch rather than failing the
run.\footnote{Implemented in \texttt{src/parser/response\_parser.py}.}
This is a deliberate robustness choice: a single malformed selection
call degrades gracefully to a specific, predictable branch rather than
discarding the entire question's data or crashing the pipeline mid-run.

As mentioned in Chapter~\ref{chap:literature}, this version of
Tree-of-Thought is simpler than the full design by Yao et
al.~\cite{yao2023tree}. It does
just one round of branching and one selection step, instead of searching
and pruning at every reasoning stage. This approach keeps the number of
calls per question fixed at four, making costs predictable, but it means
the model does not explore or backtrack within a branch after it is
created.

\section{Answer Evaluation}
\label{sec:answer-eval}

Answer correctness is determined by an automated evaluator rather than
manual grading, since the study involves 786 model responses. Before
comparison, both the ground-truth answer and the model's extracted final
answer are passed through a normalisation step that strips formatting
artefacts that would otherwise cause two mathematically identical answers
to be marked as different: code fences and bold markers, commas, dollar
signs, LaTeX inline-math delimiters, \texttt{\textbackslash text\{\}} and
\texttt{\textbackslash boxed\{\}} wrappers, and LaTeX fraction notation
(\texttt{\textbackslash frac\{a\}\{b\}}, including the brace-less
shorthand \texttt{\textbackslash frac12}) converted to plain division
(\texttt{a/b}). A literal percent sign is converted to its decimal
equivalent (e.g. ``10\%'' to $0.1$) only when the answer is purely a
number followed by \%, so that a percent sign appearing elsewhere in an
answer is not misinterpreted. Mixed numbers such as ``2 1/3'' are detected
and converted to their decimal value before general normalisation runs,
since normalisation removes whitespace and would otherwise collapse
``2 1/3'' into the digit string ``21/3''.

After normalisation, two answers are compared using three methods, and
are considered correct if any one of them agrees:

\begin{enumerate}
  \item \textbf{Numeric comparison}: both answers are converted to
  floating-point numbers (handling plain integers, decimals, and
  fractions) and compared with a small numerical tolerance, so that
  trivial floating-point representation differences are not counted as
  errors.
  \item \textbf{Symbolic comparison}: both answers are parsed as
  algebraic expressions using SymPy, and are considered equal if their
  difference simplifies to zero. This allows answers that are
  algebraically equivalent but written differently (e.g. an unsimplified
  expression versus its simplified form) to be marked correct.
  \item \textbf{Text comparison}: as a final fallback, the normalised
  answers are compared as exact strings.
\end{enumerate}

This normalisation logic was developed iteratively, whenever a review of
model outputs found a case where the correct answer was not being
extracted or matched correctly, handling for that case was added and the
check re-run.

This normalisation and comparison logic is covered by unit tests against
a set of known tricky cases (fractions, percentages, mixed numbers, and
symbolic equivalence). A separate, planned validation step -- manually
grading a sample of real model outputs by hand to check the automatic
evaluator against human judgment, rather than only hand-picked test
cases -- was not completed; this is discussed as a limitation in
Section~\ref{sec:limitations}.

\section{Research Questions and Statistical Tests}

This section explains how each research question is addressed
statistically. For the two research questions that use a formal
significance test, it covers what is being tested, the type of data
compared, and why each test is the best option instead of a simpler one.

\subsection{Chi-Square Test of Independence}

RQ1 asks whether the distribution of error types is different across the
three models. To test this, a contingency table is created with one row
for each model and one column for each of the 11 error types in the
taxonomy; ``Correct'' is a separate, twelfth classification and is not
used as a column, since only incorrect responses are classified by error
type. The table is filled with the number of incorrect responses for each
model and error type combination. The Chi-square test of independence
checks whether error type and model are unrelated, meaning each model
would have the same pattern of error types and any differences are due to
chance. Some cells in the table have small expected counts, which makes
the Chi-square approximation less reliable for those cells. This is noted
as a caution with the test result, not as a reason to dismiss it, and is
discussed more in Section~\ref{sec:limitations}.

Formally, the test statistic is
\[
\chi^2 = \sum_{i=1}^{r} \sum_{j=1}^{c} \frac{(O_{ij} - E_{ij})^2}{E_{ij}}
\]
where $O_{ij}$ is the observed count of incorrect responses for model
$i$ and error type $j$, $E_{ij} = \frac{(\text{row } i \text{
total})(\text{column } j \text{ total})}{n}$ is the expected count
under independence, $r = 3$ is the number of models, $c = 11$ is the
number of error-type categories, and $n$ is the total number of
incorrect responses. Under the null hypothesis of independence, the
statistic follows a Chi-square distribution with $(r-1)(c-1) = 20$
degrees of freedom.

\subsection{Step Count and Correctness}

RQ2 asks whether there is a descriptive relationship between step count
and answer correctness. For each model and difficulty tier, the mean and
median number of solution steps are compared between correct and
incorrect responses. Unlike RQ1 and RQ3, no formal significance test is
applied to this comparison: the analysis is exploratory, and any
differences in step count between correct and incorrect responses are
reported as observed patterns rather than statistically confirmed
effects.

\subsection{McNemar's Exact Test}

RQ3 asks whether Tree-of-Thought leads to different answer correctness
compared to Chain-of-Thought for each model. Each model answers the same
120 questions using both CoT and ToT, so the results are paired, not
independent. Comparing the raw accuracy percentages directly would
wrongly treat 240 observations as independent, when there are actually
120 pairs. McNemar's test focuses only on the discordant pairs, which are
questions where the two methods disagree: correct under CoT but not ToT
($b$), or correct under ToT but not CoT ($c$). Concordant pairs, where
both are correct or both are incorrect, do not help determine which
method is better and are left out. The exact version of the test used
here applies a binomial test to $\min(b, c)$ with $n = b + c$ trials at
$p = 0.5$. This is more suitable than the chi-square version of McNemar's
test for the sample sizes in this study.

Formally, the exact test computes a two-sided binomial p-value,
\[
p = 2 \sum_{k=0}^{\min(b,c)} \binom{n}{k} \left(\frac{1}{2}\right)^n,
\qquad n = b + c,
\]
capped at 1 if the sum exceeds 0.5, where $\binom{n}{k}$ is the number
of ways to choose $k$ successes from $n$ trials. This directly tests the
null hypothesis that a discordant pair is equally likely to favour
either prompting strategy, without relying on the chi-square
approximation McNemar's original test uses, which is unreliable at the
small $n$ values in this study ($n \leq 10$ for every model).

\subsection{Confidence Intervals for Accuracy}

Accuracy figures reported throughout Chapter~\ref{chap:results} are
accompanied by 95\% confidence intervals computed with the Wilson score
interval, rather than the normal approximation, since the normal
approximation is unreliable at small $n$ or when the observed proportion
is close to 0 or 1 -- both of which occur in this study's smaller
subgroups (Section~\ref{sec:limitations}). For $x$ correct responses out
of $n$ trials, with observed proportion $\hat{p} = x/n$ and $z$ the
standard normal critical value for the desired confidence level ($z
\approx 1.96$ for 95\%), the interval is
\[
\frac{\hat{p} + \frac{z^2}{2n} \pm z\sqrt{\frac{\hat{p}(1-\hat{p})}{n} +
\frac{z^2}{4n^2}}}{1 + \frac{z^2}{n}}.
\]

\section{Error Judging Methodology}
\label{sec:judging}

Each incorrect response is reviewed by an LLM judge: Gemini
(\texttt{gemini-flash-lite-latest}, resolving to
\texttt{gemini-3.5-flash-lite}) -- the exact same model configuration
used for one of the three evaluated models, not a separate judge-specific
variant. The judge assigns an error type from an 11-category list, notes
where the first mistake happens, and writes a brief description of the
error. It also gives a confidence level and a one-sentence explanation.

The judge is told to evaluate a model's reasoning based on its own
mathematical logic, not by comparing it to a specific solution method. As
a result, the judge only needs the correct final answer to check whether
the response is right. It does not matter that AIME questions, unlike
Hendrycks MATH, do not have official step-by-step solutions. The judge's
role is to see whether the reasoning makes sense on its own, not whether
it follows a certain method.

There is a clear conflict of interest in this judging setup. Gemini acts
as the main judge, but it is also one of the three models being tested.
This means Gemini sometimes reviews its own answers. To reduce this
risk, the judge prompt conceals the identity of the model that generated
each response. The judge is provided only with the question, category,
reference answer, and the reasoning text. Although this approach does
not completely resolve the risk of bias, as Gemini may still recognise
its own stylistic patterns, it lowers the chances of the judge favouring
responses that are easily identifiable as its own.

To address this, the evaluation included a neutral-judge validation step.
Up to 120 responses judged by Gemini were supposed to be re-judged by a
different model (\texttt{llama-3.3-70b-versatile}, via Groq). The
agreement between judges was measured using Cohen's kappa instead of just
the raw agreement rate. Raw agreement does not account for agreement that
could happen by chance, which is important because the error categories
are unevenly distributed. Cohen's kappa adjusts for this and gives a
better sense of whether the judges truly agree.

This validation step was not finished. The neutral judge model reached
Groq's free daily token limit (100,000 tokens per day) after about 20 to
25 of the planned 120 validation calls, since each call used over 4,000
tokens. Completing all 120 would have required spreading the calls over 5
or 6 days. This was a planned decision due to time limits, not a mistake,
and is reported here for transparency. More details are given as a
limitation in Section~\ref{sec:limitations}.

\section{Limitations}
\label{sec:limitations}

This section lists the main limitations of the method used in this
study. Some were mentioned earlier, but they are collected here for
completeness.

\textbf{Reasoning-effort confound (controlled, not eliminated).} As
explained in Section~\ref{sec:models}, Gemini and GPT-OSS-120B's
internal reasoning-effort settings are set to their lowest value and kept
the same for both CoT and ToT. This helps ensure that changes in internal
reasoning do not affect the comparison in RQ3. However, this only
controls the confound; it does not remove it. Setting the effort to the
minimum does not guarantee that no internal reasoning happens, especially
for GPT-OSS-120B, which may still do some reasoning even at low effort.
The exact impact on accuracy is not measured separately in this study.

\textbf{Data contamination risk.} There is some risk that the models have
seen the benchmark questions, or similar ones, during pretraining instead
of solving them through reasoning. This risk is higher for Hendrycks MATH
(used for Easy and Medium) because it is widely used for machine learning
training and evaluation. AIME (used for Hard) is also public, but the
specific AIME questions used here are comparatively recent and less
widely republished, which makes memorisation less likely, though not
impossible. This study does not eliminate either risk.

\textbf{Judge conflict of interest.} As explained in
Section~\ref{sec:judging}, Gemini is both the main error-judging model
and one of the three models being evaluated, which creates a direct
conflict of interest. The planned solution was to use a neutral judge and
measure agreement with Cohen's kappa, but this was not done because the
neutral judge's provider had a daily token limit.

\textbf{Fixed Tree-of-Thought branch count.} For each ToT question, three
candidate branches are generated (see Section~\ref{sec:prompting}). This
keeps the cost per question fixed and predictable at four model calls.
The number three was chosen for scope, not based on testing other branch
counts. Using more branches could explore more solutions per question,
but would also predictably increase costs.

\textbf{Small subgroup sample sizes.} The Hard tier categories are
unbalanced (Algebra 16, Combinatorics 16, Number Theory 11, Probability
8), so comparisons within this tier use small sample sizes and should be
viewed with caution. The Wilson score interval is used for confidence
intervals because it works better than the normal approximation with
small samples or extreme values. This concern also appears in RQ1's
Chi-square test, which was flagged due to a low minimum expected cell
count (0.29) in the Model x Error Type table.

\textbf{Pending answer-evaluator validation.} As explained in
Section~\ref{sec:answer-eval}, the automatic answer evaluator was tested with
unit tests on tricky cases. However, the planned step of manually grading
a sample of around 30--40 real model outputs, to compare the automatic
evaluator with human judgment, was not done.
